\chapter{Discussion}
\label{chap:discussion}

This chapter interprets the evaluation results from Chapter~\ref{chap:testing} in the context of
the research question, the system design choices described in Chapter~\ref{chap:implementation},
and the theoretical background established in Chapters~\ref{chap:literature} and
\ref{chap:system_design}.
The chapter begins by examining the split outcome in detail, explaining both the null kinematic
result and the positive dynamic result from first principles.
It then discusses the role of the simulator's parameter interface as the principal technical
constraint on AAS-driven fidelity improvements, reflects on what the experiment does and does not
show about AAS as a concept, and closes by revisiting the research question and acknowledging the
study's limitations.


%% ===================================================================
\section{Interpreting the Split Outcome}
\label{sec:discussion_split}
%% ===================================================================

The central finding of the experiment is that AAS-based parameter injection produced a
\textbf{split outcome}: it did not improve kinematic position accuracy, yet it produced a clear
and measurable improvement in motor current fidelity.
These two outcomes are not contradictory.
They arise from different physical mechanisms, and understanding both is necessary to evaluate what
AAS can and cannot achieve in the context of the URSim-based digital twin implemented in this
project.


%% ===================================================================
\section{Kinematic Fidelity: An Architectural Null Result}
\label{sec:discussion_kinematics}
%% ===================================================================

The joint position RMSE for Pipeline~3 (\textit{sim\_aas}) was higher than for Pipeline~2
(\textit{sim\_no\_aas}) across every joint.
At face value this appears to indicate that AAS injection made the simulation \emph{worse}.
However, the result is better understood as an architectural null result with a superimposed
timing artefact, rather than as evidence that AAS is harmful.

\subsection{Joint-Space Control and Payload Decoupling}

The pick-and-place trajectory is executed using the \texttt{moveJ} joint-space motion primitive
on the UR3 controller.
In joint-space control, the controller receives a target joint configuration and generates the
corresponding joint torques internally to drive the arm to that configuration within the specified
velocity and acceleration limits.
The payload mass parameter affects how those torques are calculated (gravity compensation and
inertial feedforward), but it does \emph{not} alter the commanded joint-angle profile.
Both Pipeline~2 and Pipeline~3 receive identical \texttt{moveJ} waypoints; the robot will
converge to the same joint angles regardless of what payload value is set.

As a consequence, no kinematic RMSE improvement through payload injection alone is
\emph{architecturally possible} within the \texttt{moveJ} control scheme.
This is not a failure of the AAS concept; it is a reflection of how the UR controller decouples
trajectory planning from dynamic compensation.
To observe a kinematic improvement from payload injection, one would need a controller mode that
exposes payload-dependent path planning, such as Cartesian force control or compliance control,
where the commanded position profile itself changes with load.

\subsection{Temporal Misalignment from Initialisation Overhead}

A second, compounding factor is the 19.5\,s initialisation overhead observed in the
\textit{sim\_aas} pipeline.
As described in Section~\ref{sec:interface_constraints}, Pipeline~3 must connect to the BaSyx
AAS server, retrieve the \textit{SimulationModel} submodel, parse the returned JSON, and apply
the parameters via RTDE before the first trajectory cycle begins.
This overhead shifts all trajectory data in Pipeline~3 forward in time relative to Pipeline~1
(\textit{real}) and Pipeline~2.

The comparison method used in \texttt{comparator.py} aligns the two runs by interpolating onto
a shared time axis from $t = 0$.
When Pipeline~3 is delayed by 19.5\,s at the start, the first comparable data point in
Pipeline~3 corresponds to a much later stage of the trajectory in Pipeline~1.
This phase shift artificially inflates the RMSE values for every joint throughout the run.
The timing artefact is therefore not a fidelity problem with the simulator itself; it is an
artefact of the evaluation methodology applied to a pipeline with a non-negligible setup latency.

These two effects together explain the null kinematic result without requiring any assumption that
AAS injection degraded performance.
The correct interpretation is: \textbf{payload injection via AAS cannot improve joint position
accuracy in a \texttt{moveJ}-controlled system, and the observed RMSE increase is an evaluation
artefact caused by initialisation-phase temporal misalignment.}


%% ===================================================================
\section{Dynamic Fidelity: AAS as an Effective Parameter Source}
\label{sec:discussion_dynamics}
%% ===================================================================

In contrast to the kinematic null result, the current measurements provide direct evidence that
AAS-based parameter injection improved simulator fidelity in a physically meaningful way.

The shoulder joint (J2) and elbow joint (J3) carry the largest gravity loads during the
pick-and-place trajectory because they support the entire distal arm mass.
Their steady-state current draw is therefore strongly coupled to the total suspended mass,
including both the arm itself and any attached tool or workpiece.
When the payload mass is set correctly to $m = 1.2$\,kg, the simulator increases the gravity
compensation and inertial feedforward torques at these joints proportionally.

The results confirm this directly.
J2 current fidelity increased from 86\,\% of real (\textit{sim\_no\_aas}) to 99\,\% of real
(\textit{sim\_aas}), an improvement of 13 percentage points.
J3 current fidelity increased from 86\,\% of real to 103\,\% of real, an improvement of
17 percentage points.
This is not a marginal statistical effect; it is a qualitative change in the simulation's
ability to reproduce the physical robot's power consumption profile at the joints that matter
most for energy and load estimation.

The overall mean RMS current rose from 0.423\,A (\textit{sim\_no\_aas}) to 0.507\,A
(\textit{sim\_aas}), narrowing the gap to the physical robot's 0.634\,A.
The remaining deficit is concentrated at J1, J5, and J6, where both simulation pipelines report
near-zero current because URSim has no friction or back-EMF model.
This residual gap is not addressable through payload injection; it requires a simulator with a
richer force-level parameter interface.

These results are consistent with the theoretical expectation for AAS-driven digital twins
described by \citeauthor{Heppner2023}~\cite{Heppner2023}, who argue that structured asset
models improve simulation fidelity precisely by providing the physical parameters that static
or manually configured simulators lack.
The experiment confirms this for the specific case of payload-dependent torque compensation in
a robot arm simulator.


%% ===================================================================
\section{The Simulator Interface as the Binding Constraint}
\label{sec:discussion_interface}
%% ===================================================================

The most significant system-level finding of this study is that the \emph{limit on AAS
contribution} is not the AAS architecture itself, but the parameter interface exposed by the
simulator.

\subsection{URSim's Narrow Parameter API}

The BaSyx AAS server used in this project stores three submodels relevant to the simulation:
\textit{MotionCommand} (payload mass, tool-centre-point offset), \textit{DynamicsParameters}
(friction coefficient, current noise level, control latency, damping factor), and
\textit{SimulationModel} (geometric and physical properties of the robot configuration).
Of these, only two scalar values from \textit{SimulationModel} are consumable by the URSim
virtual controller: the payload mass passed to \texttt{set\_payload()} and the TCP offset
passed to \texttt{set\_tcp()}.

The entire \textit{DynamicsParameters} submodel is currently unused, not because the AAS
infrastructure failed to deliver it, but because URSim provides no API endpoint to receive
friction coefficients, damping factors, or noise parameters.
This is a fundamental limitation of the simulator, not of the AAS concept.
URSim is a software replica of the UR controller designed for program validation, not for
high-fidelity physics simulation; its parameter surface was never intended to support
deep dynamic parameterisation.

\subsection{Implications for AAS Concept Validation}

This constraint has a direct consequence for how the results should be interpreted.
The experiment validates the AAS \emph{integration architecture} fully: the
digital twin successfully identified the correct AAS server at runtime, fetched the current
submodel data, parsed the structured JSON payload, and applied the parameters to the simulator
before the first trajectory cycle.
The live-fetch mechanism worked exactly as designed.

What was not validated is the broader potential of AAS-driven parameter injection, because
only one of the four available parameter categories could be injected.
A simulator with a richer API, such as a physics engine capable of accepting friction
coefficients and joint-level damping parameters (for example, a Gazebo model or a
MATLAB Simscape model), would be able to consume the full \textit{DynamicsParameters}
submodel.
Based on the J2 and J3 current results, there is strong reason to expect that such a richer
injection would also improve the current fidelity at J1, J5, and J6, where friction-related
discrepancies currently dominate.


%% ===================================================================
\section{Comparison with Related Work}
\label{sec:discussion_related}
%% ===================================================================

The study contributes a concrete implementation and quantitative evaluation to a research area
that, until recently, has been characterised primarily by architectural proposals and
demonstrator systems without systematic empirical comparison.

\citeauthor{Bu2025}~\cite{Bu2025} and \citeauthor{Zhang2025}~\cite{Zhang2025} both propose AAS
as a parameter source for simulation, but focus on discrete-event and process-level simulation
rather than real-time kinematic and dynamic fidelity in robot motion control.
The present study extends the scope to the joint-level physics of a real industrial arm, where
the physics of gravity compensation provide a measurable and directly interpretable signal.

The approach of using a physical robot as the ground-truth reference and quantifying per-joint
RMSE and RMS current aligns with the validation methodology recommended by
\citeauthor{vuzem2025dt}~\cite{vuzem2025dt} for automated simulation model generation.
The quasi-experimental design with paired comparisons (Pipeline~1 versus Pipeline~2 and
Pipeline~1 versus Pipeline~3) follows the structure advocated by
\citeauthor{cook1979quasi}~\cite{cook1979quasi} for controlled comparisons in settings where
full randomisation is not feasible.

The finding that AAS injection improves dynamic fidelity but not kinematic fidelity under
joint-space control is, to the authors' knowledge, not previously reported.
It highlights a category distinction that future AAS-simulation integration studies should
address explicitly: the distinction between \emph{kinematic fidelity} (position-level accuracy)
and \emph{dynamic fidelity} (torque, current, and energy accuracy), which respond differently to
mass and friction parameters in a joint-space-controlled system.


%% ===================================================================
\section{Revisiting the Research Question}
\label{sec:discussion_rq}
%% ===================================================================

The research question posed in Chapter~\ref{chap:problem_formulation} asks how AAS data
structures for a UR3 robotic arm can lead to more accurate and consistent simulation results
compared with conventional, manually parameterised approaches.

The experiment provides a two-part answer.

\textbf{In terms of dynamic fidelity}, the AAS approach demonstrated a clear advantage.
By providing a machine-readable, live-updated source of the correct payload mass, the AAS
enabled Pipeline~3 to reproduce the physical robot's joint current profile at the primary
load-bearing joints (J2 and J3) with high accuracy: 99\,\% and 103\,\% of real values,
respectively, compared with only 86\,\% for the manually parameterised baseline.
This is precisely the intended function of the AAS in the Industry~4.0 context: not just storing
parameters, but making them available as an authoritative, structured source for connected
systems~\cite{IDTA_AAS_Part1, Bader2021}.

\textbf{In terms of kinematic fidelity}, the AAS approach showed no advantage in this
experimental configuration.
The joint-space control architecture decouples payload from trajectory execution, and the
initialisation overhead introduced a temporal misalignment that dominated the RMSE comparison.
These are specific properties of the chosen simulator and control mode, not general limitations
of AAS-driven digital twins.

The answer to the research question is therefore \emph{conditionally affirmative}: AAS parameter
injection does improve simulation accuracy at the dynamic level, and the infrastructure required
to do so is fully realisable with current technology.
The benefit is, however, bounded by the parameter surface of the simulator, and is not
observable at the kinematic level under \texttt{moveJ} joint-space control.


%% ===================================================================
\section{Limitations}
\label{sec:discussion_limitations}
%% ===================================================================

Several limitations constrain the generalisability of the findings and should be considered when
interpreting the results.

\textbf{Single experiment run per condition.}
Pipeline~2 and Pipeline~3 were each evaluated using a single representative run (\textit{db4c6d73}
and \textit{7b8c4a8e}, respectively).
While the physical pipeline was repeated six times to provide variability estimates for the
reference, the simulation pipelines were not.
Timing jitter and non-deterministic RTDE latency in the simulation environment mean that
a single run may not fully represent the distribution of possible outcomes.

\textbf{Limited trajectory complexity.}
The evaluation trajectory is a three-cycle pick-and-place motion with a fixed set of waypoints.
More complex trajectories with higher accelerations, orientation changes, and payload dynamics
would stress-test the AAS injection across a wider operating envelope.

\textbf{One simulator platform.}
All simulation runs use URSim, whose narrow parameter API is the principal constraint on AAS
utilisation.
The findings cannot be directly extrapolated to physics engines (such as Gazebo or
MATLAB Simscape) that expose richer parameter surfaces.

\textbf{Single payload value.}
Only one AAS-injected payload mass was evaluated ($m = 1.2$\,kg).
A systematic sweep across payload values would allow the relationship between mass, RMSE, and
current fidelity to be characterised quantitatively and would reveal whether the accuracy
improvements scale predictably.

\textbf{No submodel update during execution.}
The AAS parameters were fetched once before the run and held constant.
A fully dynamic scenario, in which the submodels are updated mid-run to reflect a changing
workpiece mass, was not tested.
Such a scenario would demonstrate the live-update capability that is one of the stated
advantages of the AAS approach~\cite{IDTA_AAS_Part1}.

\textbf{Delimitations from scope.}
As established in Section~\ref{sec:delimitations}, this study deliberately focuses on a single
UR3 robot executing a single pick-and-place task.
Safety systems, multi-robot coordination, process-level scheduling, and the long-term reliability
of the AAS server infrastructure are all outside scope.
